% -----------------------------------------------------------------------------------------------------
\section{多腿振幅的红外极点}\label{sec:sudakov-multileg}

$n$ 条有色外线的在壳振幅同样有至多 $1/\varepsilon^{2L}$ 的红外极点，且由于软胶子可以连接任意两条外线，极点带有色关联。自然的语言是\emph{色空间形式}：
振幅是色空间矢量 $|\mathcal M(\{p\},\varepsilon)\rangle$，$\mathbf T_i$ 作用于第 $i$ 个部分子的色指标（出射夸克为 $t^a$，出射反夸克为 $-t^{a\,\tp}$，胶子为伴随表示生成元，入射粒子取共轭），
$\mathbf T_i^2=C_i$，色守恒为 $\sum_i\mathbf T_i|\mathcal M\rangle=0$。形状因子的结构推广为色空间中的乘性重整化因子 $\mathbf Z$，两圈以内其反常维度只含两两的“偶极子”色关联。
$\mathbf Z$ 的结构来自 $n$ 条 Wilson 线的软关联函数，我们先说明这种关联函数如何指数化。

\paragraph{非阿贝尔指数化。}非阿贝尔理论中 Wilson 线的色矩阵按路径排序，指数化仍成立，但指数中的图（web）带修正的色因子：例如两线间的两胶子交换，
不交叉梯形图（色因子 $C_F^2$）被单圈项的平方完全重现，交叉梯形图（$C_F(C_F-\frac{C_A}2)$）只留下 $-\frac12C_FC_A$。% check: s4_wilson.py (h)

\begin{proposition}[非阿贝尔指数化（webs） \tagA]\label{prop:sudakov-webs}
(a) 在原点相遇的 Wilson 线的真空关联函数可以写成 $\exp[\sum_w\tilde C(w)\,\mathcal F(w)]$，求和遍及 web（对 Wilson 线的“不可约”图），$\tilde C(w)$ 是修正的色因子；
对两条线，web 的修正色因子由“最大非阿贝尔”部分组成\cite{Gatheral1983,FrenkelTaylor1984}（附录~B 的\cref{prop:appB-webs}）。
(b)（$O(g^2)$ 内容，\tagR，$R_{\mathrm{def}}$）设 $W[0,-\infty n]$ 处于由有限个平面波组成的背景场中，第 $r$ 个模式的顶点因子记为矩阵 $T_r$（含 $\ii g\,n\cdot\tilde a_r$ 与色生成元），
$E(r)=\frac{\ii}{n\cdot k_r+\ii\epsilon}$，$E(r,s)=\frac{\ii}{n\cdot k_r+\ii\epsilon}\frac{\ii}{n\cdot(k_r+k_s)+2\ii\epsilon}$（绝热正规子，\cref{prop:prelim-wilson-rules}），
则 $W^{(1)}=\sum_rT_rE(r)$、$W^{(2)}=\sum_{r,s}T_sT_rE(r,s)$，并且在有限 $\epsilon$ 下\emph{严格}成立
\begin{equation}\label{eq:sudakov-web2}
W^{(2)}-\tfrac12\big(W^{(1)}\big)^2=\tfrac14\sum_{r,s}[T_s,T_r]\,\big(E(r,s)-E(s,r)\big):
% VF-WL-3
\end{equation}
$\ln W=W^{(1)}+\big[W^{(2)}-\frac12(W^{(1)})^2\big]+O(g^3)$ 的 $O(g^2)$ 部分只含对易子；阿贝尔情形 $\ln W=W^{(1)}$。（本命题只涉及 Wilson 线算符。）
\uses{def:prelim-wilson,prop:prelim-wilson-rules,lem:prelim-eikonal-identity,prop:appB-webs}
\end{proposition}

\begin{proofsketch}
(a) 是 Gatheral 以及 Frenkel 与 Taylor 的组合学论证（对色因子按“最大非阿贝尔”部分重新组织），其概要见\cref{prop:appB-webs}。
(b) 的完整证明：$W^{(1)}$、$W^{(2)}$ 是\cref{prop:prelim-wilson-rules}的一阶与二阶项（离 $-\infty$ 近的模式 $r$ 在右）。由程函恒等式在有限 $\epsilon$ 下的形式（\cref{lem:prelim-eikonal-identity}，$m=2$），
$E(r)E(s)=E(r,s)+E(s,r)$。于是 $\frac12(W^{(1)})^2=\frac12\sum_{r,s}T_sT_r[E(r,s)+E(s,r)]$，
$W^{(2)}-\frac12(W^{(1)})^2=\frac12\sum_{r,s}T_sT_r[E(r,s)-E(s,r)]$；在第二项中交换求和指标 $r\leftrightarrow s$ 得 $\frac12\sum_{r,s}[T_s,T_r]E(r,s)$，再对 $(r,s)$ 反对称化即得式~\eqref{eq:sudakov-web2}。
若 $T_r$ 彼此对易，右边为零。
\end{proofsketch}

\paragraph{地位与缺口}(a) 是全阶的组合学论证（[A]），对在一点相遇的两条线在文献中完整给出\cite{Gatheral1983,FrenkelTaylor1984}；多条线的 web 及其重整化本书不需要。
(b) 是严格的代数恒等式（[R]）；它依赖于绝热正规子——若每个程函分母都加同一个 $+\ii\epsilon$，程函恒等式在有限 $\epsilon$ 下不再成立（附录~B 的\cref{lem:appB-shuffle}），式~\eqref{eq:sudakov-web2} 只在 $\epsilon\to0$ 的分布意义下成立。

\begin{theorem}[无质量多腿振幅红外极点的因子化 \tagA]\label{thm:sudakov-ir-poles}
设 $|\mathcal M(\{p\},\varepsilon)\rangle$ 是 $n$ 个无质量部分子的紫外已重整化的在壳振幅（红外由维数正规化正规化）。则存在色空间中的矩阵 $\mathbf Z(\{p\},\mu,\varepsilon)$（只含 $\varepsilon$ 的极点），
使 $|\mathcal M_{\rm fin}(\{p\},\mu)\rangle\equiv\mathbf Z^{-1}|\mathcal M(\{p\},\varepsilon)\rangle$ 在 $\varepsilon\to0$ 时有限：
\begin{equation}\label{eq:sudakov-Zfact}
|\mathcal M(\{p\},\varepsilon)\rangle=\mathbf Z(\{p\},\mu,\varepsilon)\,|\mathcal M_{\rm fin}(\{p\},\mu)\rangle,\qquad
\boldsymbol\Gamma=-\mathbf Z^{-1}\frac{\dd\mathbf Z}{\dd\ln\mu},\qquad \frac{\dd}{\dd\ln\mu}|\mathcal M_{\rm fin}\rangle=\boldsymbol\Gamma\,|\mathcal M_{\rm fin}\rangle .
\end{equation}
两圈以内，反常维度取偶极子形式
\begin{equation}\label{eq:sudakov-dipole}
\boldsymbol\Gamma=\sum_{(i,j)}\frac{\mathbf T_i\cdot\mathbf T_j}{2}\,\gamma_{\rm cusp}(\alpha_s)\ln\frac{\mu^2}{-s_{ij}}+\sum_i\gamma^i(\alpha_s),\qquad s_{ij}=2\sigma_{ij}\,p_i\cdot p_j+\ii0,
% VF-SUD-5
\end{equation}
其中 $(i,j)$ 遍历 $i\ne j$ 的全部\emph{有序}对（每个无序对计两次），$\sigma_{ij}=+1$（$i,j$ 同为入射或同为出射）或 $-1$（其余），
$\gamma_{\rm cusp}=4\frac{\alpha_s}{4\pi}+O(\alpha_s^2)$ 是去掉 Casimir 的 cusp 反常维度（$\Gamma^F_{\rm cusp}=C_F\gamma_{\rm cusp}$，\cref{def:sudakov-cusp}），$\gamma^q=-3C_F\frac{\alpha_s}{4\pi}+O(\alpha_s^2)$。
在单圈阶，$\ln\mathbf Z=\frac{\alpha_s}{4\pi}\big[\frac{\boldsymbol\Gamma'_0}{4\varepsilon^2}+\frac{\boldsymbol\Gamma_0}{2\varepsilon}\big]+O(\alpha_s^2)$，其中 $\boldsymbol\Gamma=\frac{\alpha_s}{4\pi}\boldsymbol\Gamma_0+\cdots$，
$\boldsymbol\Gamma'\equiv\partial\boldsymbol\Gamma/\partial\ln\mu=-\gamma_{\rm cusp}\sum_iC_i$。对夸克形状因子（$n=2$），式~\eqref{eq:sudakov-dipole} 重现硬匹配系数的演化 $\dd\ln C_V/\dd\ln\mu$（见下例）。（本定理只涉及部分子振幅。）
\uses{thm:sudakov-factorization,def:sudakov-cusp,prop:sudakov-webs,prop:sudakov-oneloop,ex:sudakov-onshell-hard}
\end{theorem}

\begin{proofsketch}
(1) 把\cref{thm:sudakov-factorization}的论证推广到 $n$ 条宽角外线，得到振幅的硬—喷注—软因子化\cite{Akhoury1979,Sen1983}（全阶的振幅表述见\cite{FeigeSchwartz2014,Ma2020}）。
(2) 纯维数正规化中喷注与软矩阵元无标度，其红外与紫外极点相消（\cref{ex:sudakov-onshell-hard}的机制），振幅的红外极点由低能矩阵元的紫外重整化因子给出；喷注部分色对角，
软部分是 $n$ 条类光 Wilson 线的关联函数，在色空间中乘性重整化，web 结构（\cref{prop:sudakov-webs}）保证它是矩阵的指数\cite{Becher:2009cu,Becher:2009qa,DixonMagneaSterman2008}。
(3) 于是 $\mathbf Z$ 满足式~\eqref{eq:sudakov-Zfact}。设 $\mathbf Z=1+\frac{\alpha_s}{4\pi}(\frac{\mathbf z_2}{\varepsilon^2}+\frac{\mathbf z_1}{\varepsilon})+O(\alpha_s^2)$，由 $\frac{\dd\alpha_s}{\dd\ln\mu}=-2\varepsilon\alpha_s+O(\alpha_s^2)$，
$\frac{\dd\mathbf Z}{\dd\ln\mu}=\frac{\alpha_s}{4\pi}[-\frac{2\mathbf z_2}{\varepsilon}-2\mathbf z_1+\frac1\varepsilon\frac{\partial\mathbf z_1}{\partial\ln\mu}]$；在 $O(\alpha_s)$ 令其等于有限的 $-\mathbf Z\boldsymbol\Gamma=-\frac{\alpha_s}{4\pi}\boldsymbol\Gamma_0+O(\alpha_s^2)$，得 $\mathbf z_1=\frac12\boldsymbol\Gamma_0$、$\mathbf z_2=\frac14\boldsymbol\Gamma'_0$。
由色守恒 $\sum_{(i,j)}\mathbf T_i\cdot\mathbf T_j=-\sum_iC_i$，故 $\boldsymbol\Gamma'=-\gamma_{\rm cusp}\sum_iC_i$ 色对角。
(4) 单圈时软关联函数是两两 Wilson 线之间单胶子交换之和，每对贡献色因子 $\mathbf T_i\cdot\mathbf T_j$ 与类光尖点反常维度（\cref{def:sudakov-cusp}）；两圈时软反常维度矩阵与单圈矩阵成正比\cite{Aybat:2006mz}。
Becher 与 Neubert\cite[式~(7)]{Becher:2009cu}（亦见\cite[式~(7)]{Becher:2009qa}）以及 Gardi 与 Magnea\cite{Gardi:2009qi}猜测它对所有阶成立；两圈的显式结构与 Catani 公式\cite{Catani:1998bh}及\cite{Sterman:2002qn}一致。
\end{proofsketch}

\paragraph{地位与缺口}极点的全阶因子化（式~\eqref{eq:sudakov-Zfact}）是全阶论证（[A]），它继承振幅层次硬—喷注—软因子化的地位（与\cref{thm:sudakov-factorization}同样的两个瓶颈：
减除的 E-S 地位与非阿贝尔 Ward 恒等式）。偶极子形式在单圈与两圈是显式计算的结果（[R] 意义上的定阶内容，两圈见\cite{Aybat:2006mz}）；
作为全阶精确公式它不成立（\cref{prop:sudakov-dipole-violation}）。升级为 $[R]$ 需要振幅因子化的完整证明（参见第~24~章\cref{prop:front-ir-forest}的地位，前向指针）。

\begin{example}[两部分子情形与色守恒]\label{ex:sudakov-dipole-twoparton}
$n=2$ 时色单态要求 $\mathbf T_1+\mathbf T_2=0$，$\mathbf T_1\cdot\mathbf T_2=-C_F$，两个有序对给出 $\boldsymbol\Gamma=-C_F\gamma_{\rm cusp}\ln\frac{\mu^2}{-s_{12}}+2\gamma^q$。类空形状因子一入一出，
$\sigma_{12}=-1$，$-s_{12}=Q^2$，故 $\boldsymbol\Gamma=\frac{C_F\alpha_s}{4\pi}(4L-6)$（$L=\ln\frac{Q^2}{\mu^2}$）；而式~\eqref{eq:sudakov-F1-L} 的有限部分 $C_V=1+\frac{C_F\alpha_s}{4\pi}(-L^2+3L-8+\zeta_2)$
满足 $\frac{\dd C_V}{\dd\ln\mu}=\frac{C_F\alpha_s}{4\pi}(-2L+3)(-2)$，两者一致。把 $(i,j)$ 读作无序对则少一半——这就是“有序对”的含义（Becher 与 Neubert\cite{Becher:2009qa}称同一约定为“unordered tuples”，即 $(i,j)$ 与 $(j,i)$ 都计入）。% VF-SCET-8, VF-SUD-5
极点：$\boldsymbol\Gamma_0=C_F(4L-6)$、$\boldsymbol\Gamma'_0=-8C_F$，单圈公式给出 $\ln\mathbf Z=\frac{C_F\alpha_s}{4\pi}[-\frac2{\varepsilon^2}+\frac{2L-3}\varepsilon]$，正是式~\eqref{eq:sudakov-F1-L} 的极点部分；
类时形状因子 $\sigma_{12}=+1$，$\ln\frac{\mu^2}{-s-\ii0}=\ln\frac{\mu^2}{s}+\ii\pi$，$\ln\mathbf Z$ 获得虚部 $\frac{C_F\alpha_s}{4\pi}(-\frac{2\pi\ii}\varepsilon)$（与后面的\cref{prop:sudakov-coulomb}一致）。
对 $q\bar qg$，$\sum_{(i,j)}\mathbf T_i\cdot\mathbf T_j=-(2C_F+C_A)$。% check: s3_irpoles.py (a)-(c)
\uses{thm:sudakov-ir-poles,prop:sudakov-oneloop}
\end{example}

%%% BEGIN fig:sudakov-dipole
\begin{figure}[htbp]
\centering
\begin{tikzpicture}[fd,
  fa soft/.style={gluon,green!45!black},
  fa web/.style={soft blob,line width=0.7pt,pattern=north east lines,pattern color=green!45!black},
  fa over/.style={preaction={draw=white,line width=4.2pt,decorate=false}},
  fa lab/.style={font=\small},
  fa note/.style={font=\footnotesize},
  fa cap/.style={font=\footnotesize},
  fa legs/.pic={\foreach \a/\n in {66/1,22/2,-22/3,-66/4} {
                  \draw[plain] (\a:0.4) -- (\a:2.2) node[font=\small,anchor=\a+180,inner sep=2pt] {$\n$};}
                \node[hard blob,circle,minimum size=0.8cm,inner sep=0pt] at (0,0) {$H$};}]
  % ===================== (a) dipole terms =====================
  \begin{scope}
    \pic at (0,0) {fa legs};
    \draw[fa soft] (66:1.75) to[bend left=38] (22:1.75);
    \draw[fa soft] (-22:1.75) to[bend left=38] (-66:1.75);
    \foreach \a in {66,22,-22,-66} \node[vertex] at (\a:1.75) {};
    \node[fa note,anchor=south west,inner sep=1pt] at (44:2.12) {$\mathbf T_1\cdot\mathbf T_2$};
    \node[fa note,anchor=north west,inner sep=1pt] at (-44:2.12) {$\mathbf T_3\cdot\mathbf T_4$};
    \node[fa note,align=left,anchor=west] at (2.75,0) {对所有\\$i\neq j$ 求和};
    \node[fa cap] at (1.3,-3.05) {(a)\ 两圈以内：偶极子（色关联 $\mathbf T_i\cdot\mathbf T_j$）};
  \end{scope}
  % ===================== (b) three-loop webs (non-dipole) =====================
  % web blob in the free sector between legs 2 and 3; gluons that pass over a leg get a white gap
  \begin{scope}[xshift=7.2cm]
    \pic at (0,0) {fa legs};
    \draw[fa soft,fa over] (1.66,0.27) -- (66:1.5);
    \draw[fa soft] (1.86,0.2) -- (22:2.02);
    \draw[fa soft] (1.86,-0.2) -- (-22:2.02);
    \node[fa web,circle,minimum size=0.62cm,inner sep=0pt] (W) at (1.75,0) {};
    \foreach \p in {(66:1.5),(22:2.02),(-22:2.02)} \node[vertex] at \p {};
    \node[fa note] at (0.9,-2.6) {连接三条外线};
  \end{scope}
  \begin{scope}[xshift=11.6cm]
    \pic at (0,0) {fa legs};
    \draw[fa soft,fa over] (1.66,0.27) -- (66:1.5);
    \draw[fa soft,fa over] (1.66,-0.27) -- (-66:1.5);
    \draw[fa soft] (1.86,0.2) -- (22:2.02);
    \draw[fa soft] (1.86,-0.2) -- (-22:2.02);
    \node[fa web,circle,minimum size=0.62cm,inner sep=0pt] (W) at (1.75,0) {};
    \foreach \p in {(66:1.5),(22:2.02),(-22:2.02),(-66:1.5)} \node[vertex] at \p {};
    \node[fa note] at (0.9,-2.6) {连接四条外线};
  \end{scope}
  \node[fa cap] at (9.45,-3.05) {(b)\ 三圈起：非偶极子 web（阴影）};
  % ===================== anomalous dimension (VF-SUD-5) =====================
  \node[font=\small] at (5.6,-3.9)
    {$\boldsymbol\Gamma=\displaystyle\sum_{(i,j)}\frac{\mathbf T_i\cdot\mathbf T_j}{2}\,\gamma_{\rm cusp}(\alpha_s)\ln\frac{\mu^2}{-s_{ij}}+\sum_i\gamma^i(\alpha_s)$
     \quad（两圈以内成立）};
\end{tikzpicture}
\caption{多腿振幅的红外结构：(a) 两圈以内由偶极子项（色关联 $\mathbf T_i\cdot\mathbf T_j$）给出，反常维度如图下公式（VF-SUD-5，$(i,j)$ 取有序对）；
(b) 三圈起出现连接三条或四条外线的非偶极子 web（偶极子公式作为全阶公式在三圈失效，[V]）。}
\label{fig:sudakov-dipole}
\end{figure}
%%% END fig:sudakov-dipole

\cref{fig:sudakov-dipole}(a) 是两圈以内的偶极子结构；从三圈起一个 web 可以同时连接三条或四条 Wilson 线（\cref{fig:sudakov-dipole}(b)），带来新的色张量与运动学依赖。

\begin{proposition}[偶极子公式作为全阶精确公式的失效 \tagV]\label{prop:sudakov-dipole-violation}
式~\eqref{eq:sudakov-dipole} 作为对所有阶成立的公式是错误的：三圈时 $\boldsymbol\Gamma$ 获得非偶极子修正 $\boldsymbol\Delta^{(3)}_n$。
它由两类色结构组成：只含三个不同部分子的项 $\propto C\,f^{abe}f^{cde}\{\mathbf T_i^a,\mathbf T_i^d\}\mathbf T_j^b\mathbf T_k^c$，$C=\zeta_5+2\zeta_2\zeta_3\approx4.99$，对三个有色部分子是一个与运动学无关的常数；
含四个不同部分子的项乘以共形交比的权五单值多重对数函数，对四个及以上有色部分子依赖于运动学。该修正保持类时严格共线因子化\cite[式~(7)、(13)、(14)、(17)]{Almelid:2015jia}。
（本命题只涉及部分子振幅。）
\uses{thm:sudakov-ir-poles}
\end{proposition}

\begin{proofsketch}
（证据：显式计算。）Almelid、Duhr 与 Gardi 计算了连接三条与四条 Wilson 线的全部三圈 web（先取非类光的线以避免共线奇异性，再取类光极限）\cite{Almelid:2015jia}；Henn 与 Mistlberger 在 $\mathcal N=4$ 超对称 Yang--Mills 理论中的三圈四胶子振幅
为此提供了独立检验\cite{Henn:2016jdu}；此前 Caron-Huot 由 Regge 极限中的 BFKL 动力学论证了偶极子猜测从四圈起不成立\cite[第~4.3~节，式~(4.28)]{Caron-Huot:2013fea}。
结构上：三线与四线 web 的色因子含两个 $f^{abc}$ 的缩并，不能约化为 $\mathbf T_i\cdot\mathbf T_j$；类光 Wilson 线在 $n_i\to\kappa_in_i$ 下的重标度不变性要求运动学依赖只通过共形交比进入，
三个部分子构不成共形交比，故只剩常数。
\end{proofsketch}

\paragraph{地位与缺口}[V]，证据类型：显式三圈计算（\cite{Almelid:2015jia}，并由\cite{Henn:2016jdu}独立检验）以及四圈的 Regge 论证\cite{Caron-Huot:2013fea}。
被推翻的只是“偶极子形式对所有阶精确”这一陈述；\cref{thm:sudakov-ir-poles}的极点因子化本身（[A]）与两圈以内的偶极子形式不受影响。
三圈修正首先出现在三圈振幅的 $1/\varepsilon$ 极点中；它与类时共线极限相容（严格共线因子化得以保持），这是一个非平凡的检验。
